% Generated draft: requires mathematical and source review.
\documentclass[11pt]{article}
\usepackage[margin=25mm]{geometry}
\usepackage{fontspec}
\setmainfont{Latin Modern Roman}
\usepackage{amsmath,amssymb,mathtools,unicode-math}
\setmathfont{Latin Modern Math}
\usepackage{enumitem,needspace,xurl,hyperref}
\setlength{\parindent}{0pt}
\setlength{\parskip}{5pt}
\setlength{\emergencystretch}{3em}
\newcommand{\N}{\mathbb{N}}
\newcommand{\Z}{\mathbb{Z}}
\newcommand{\Q}{\mathbb{Q}}
\newcommand{\R}{\mathbb{R}}
\newcommand{\C}{\mathbb{C}}
\newcommand{\F}{\mathbb{F}}
\newcommand{\A}{\mathbb{A}}
\newcommand{\PP}{\mathbb{P}}
\newcommand{\E}{\mathbb{E}}
\newcommand{\eps}{\varepsilon}
\newcommand{\dd}{\,\mathrm{d}}
\begin{document}
\section*{UnsolvedMath ID: 30000147}
\textbf{Source problem code:} catalogue entry
\subsection{Definitions and mathematical statement}
The source asks whether every nondegenerate polynomial phase has a power-decay estimate for its associated oscillatory integral. In general notation, an oscillatory integral has the form
\[
I_\lambda(P;F)=\int_{\R^m}e^{i\lambda P(x)}A_F(x)\,\dd x,
\]
where \(P\) is real-valued, \(\lambda\) is the large oscillation parameter, and \(A_F\) ranges over the source's intended admissible amplitudes. Power decay means an estimate of the form
\[
|I_\lambda(P;F)|\leq C_P(1+|\lambda|)^{-\varepsilon_P}\mathcal N(F),
\qquad\varepsilon_P>0,
\]
for all admissible inputs and large \(|\lambda|\), with constants independent of \(\lambda\) and of inputs normalized by \(\mathcal N\). The source does not specify the input class, normalization, or permitted dependence of the constants; it does not state a uniform exponent over all phases.

For example, in a multilinear projection setting one would have
\[
A_F(x)=\eta(x)\prod_{j=1}^r f_j(\pi_jx),
\qquad\mathcal N(F)=\prod_j\|f_j\|_\infty,
\]
and nondegeneracy would typically be relative to the maps \(\pi_j\), rather than mean nonsingularity of the Hessian of \(P\). This is an illustration of the missing structural data, not a claim that these exact maps, norms, or definition are supplied in the record. The original integral, support and smoothness conditions on its cutoff, dimension, and definition of nondegenerate must be recovered before the displayed template becomes a definite conjecture. Replacing the problem by a scalar stationary-phase statement for smooth fixed amplitudes could change its content substantially.

The supplied literature assessment checked 2026-08-21 reports that the general power-decay question remains open in the cited record, even for quadratic phases in three variables. Because its nondegeneracy and integral conventions are not reproduced, that reported unresolved case must not be read as a claim about every ordinary quadratic oscillatory integral.
\subsection{Short English statement}
Ask whether the intended nondegeneracy condition on a polynomial phase guarantees some positive power of decay in the oscillation parameter. The integral, amplitude class, norms, and nondegeneracy definition are missing and cannot be replaced by a different stationary-phase problem.
\subsection{Sources}
\begin{itemize}
\item Local UnsolvedMath registry, permanent ID 30000147, catalogue entry, \emph{Power Decay for Polynomial Oscillatory Integrals}.
\item \emph{Nonlinear Waves and Dispersive Equations} (2005), Oberwolfach Reports, pp. 2653--2728, DOI \url{https://doi.org/10.4171/owr/2004/50}.
\item Supplied reference labeled ``current EMS report text'': \url{https://ems.press/content/serial-article-files/45969}.
\item Supplied literature triage checked 2026-08-21 and statement-verification claim of comparison with the official report. No explicit oscillatory-integral formula or nondegeneracy definition is reproduced.
\end{itemize}
\end{document}
